% Chapter 12: GF(2)-Linear Retrieval
% DEFINES: ch:gf2-retrieval; sec:retrieval-no-gate, sec:designable-junk,
%   sec:rank-threshold, sec:structural-hiding, sec:ch12-notes;
%   thm:span-output, rem:rank-threshold, rem:freq-hiding, tab:hiding
% REFERENCES (backward): def:totality, sec:acceptance, tab:levers,
%   sec:compositional-measures, def:uniformity
% REFERENCES (forward, part labels): part:frontiers

\chapter{GF(2)-Linear Retrieval}
\label{ch:gf2-retrieval}

The hash-based construction of \Cref{ch:hash-constructions} hides frequency
the way \Cref{ch:marginal-scale} prescribed: by spending tokens, more
representations for frequent values, a cost paid per element. This chapter is
a construction of a different character, where the frequency-hiding is not
bought but \emph{structural}, a property of a linear system that holds at
zero per-query cost. It is summarized from the codec paper, a standalone
submission; the constructions and the threshold theorem are stated with their
intuition and pointed to the paper for proof.\footnote{The construction, the
rank threshold, and the frequency-hiding game are \cite{towell2026codec},
which also carries the header-only C++23 implementation and the at-scale
empirical checks; this chapter states and motivates them.}

\section{Retrieval Without a Membership Gate}
\label{sec:retrieval-no-gate}

A retrieval data structure, a ribbon, XOR, or Bloomier filter, returns a
value for \emph{every} key, with no membership test: the stored value for a
member, and for a non-member some value the literature dismisses as junk, a
don't-care. This is totality (\Cref{def:totality}) in a new guise: every
input produces output, and there is no out-of-band signal that a key was
never stored. The structures of interest here are linear. In a homogeneous
ribbon retrieval over a key set $S$, a query is answered by a GF(2)-linear
combination of the structure's stored rows, so the value returned is an XOR
of solution vectors selected by the key's hash. The whole machinery is a
linear system over $\mathrm{GF}(2)$, and that linearity is what makes the junk
analyzable.

\section{Junk Is Designable}
\label{sec:designable-junk}

The dismissed value, the junk, is the chapter's object, and the result is
that it is not arbitrary at all.

\begin{theorem}[Non-member output, summarized]
\label{thm:span-output}
For homogeneous ribbon retrieval, the value returned for a non-member key is
uniform on the GF(2) span $W$ of the stored codewords. Its distribution is
therefore fixed by a \emph{public} value codec, the map from values to
codewords, and is provably independent of how often each value was stored.
\end{theorem}

The consequence is that junk is a \emph{designable} object. The non-member
output distribution does not leak the stored frequencies, because it depends
only on the span $W$ and the public codec, not on the private data; a
designer chooses the codec and thereby chooses what non-members reveal,
which can be nothing about frequency at all. This is the acceptance predicate
of \Cref{sec:acceptance} realized at the level of a linear code: the codec
plays the role of the acceptance partition, and the span plays the role of
the populated support.

\section{The Rank Threshold}
\label{sec:rank-threshold}

Whether the codec achieves full control over the non-member distribution is
not a matter of degree but of a sharp threshold.

\begin{remark}[The rank threshold]
\label{rem:rank-threshold}
For a balanced codec with $K$ classes, full control holds if and only if the
stored codewords are transversal to the codec's class partition, equivalently
\[
  \rank\bigl(\proj|_{W}\bigr) \;=\; \log_2 K,
\]
the projection of the stored span onto the codec's class coordinates having
full rank. This is a \emph{step} in an integer rank, with no graded middle:
above the threshold the codec controls every class, and below it whole codec
classes receive exactly zero probability. A second invariant, the matroid
cogirth $d^{*}$ of the stored column system, is the exact adversarial erasure
budget before control collapses, at least $\tfrac{K}{2}\,m_{\min}$ under full
support. Rank governs whether control holds; cogirth governs how robust that
control is to erasures.
\end{remark}

The step is worth dwelling on, because it is unlike the smooth $\delta$ of
\Cref{def:uniformity}. There is no codec that ``mostly'' controls the
output: either the rank condition holds and the public codec fixes the
distribution exactly, or it fails and a hard zero appears in some class. The
designer's task is a discrete one, ensure the rank, then choose the codec,
rather than a continuous tuning.

\section{Frequency-Hiding at Zero Per-Query Cost}
\label{sec:structural-hiding}

The payoff is a confidentiality property the hash-based construction cannot
match in kind. Consider a distinguishing game in which an adversary, with
oracle access to non-member outputs, tries to tell a high-frequency stored
value from a low-frequency one. In the idealized model the two non-member
laws are \emph{identical}, both uniform on the span, so the adversary's
advantage is exactly zero, for any number of adaptive queries; there is
nothing to learn because the output law does not depend on frequency at all.
In the real construction the two laws lie within total variation
$\delta(p_0) + \delta(p_1)$, a small, scale-independent, redundancy-governed
term, so a single observation distinguishes them by at most that distance and
more queries only sharpen an estimate of two fixed laws whose separation does
not grow with the frequency gap.

The decisive word is structural. The hiding here is a property of the linear
system, established once when the structure is built, and it costs
\emph{nothing per query} (\Cref{tab:hiding}). This is the contrast with the
homophonic multiplicity of \Cref{tab:levers}, which hides frequency by
spending tokens and which an online defense pays for on every access. The two
are genuinely different mechanisms for the same end, and they pull in
different directions under a second adversary. The multi-instance coincidence
oracle of \Cref{sec:compositional-measures} is blunted not by a flat codec
but by a \emph{concentrated} one, the opposite of what the frequency-hiding
game wants, so there is no single optimal codec, only a public tunable
frontier between the two defenses.

\begin{table}[t]
\centering
\caption{Two ways to hide frequency.}
\label{tab:hiding}
\begin{tabular}{@{}lll@{}}
\toprule
Mechanism & How & Cost \\
\midrule
Homophonic multiplicity (\Cref{tab:levers}) & spread frequent values over tokens & space; per-query for online defenses \\
GF(2) codec (\Cref{thm:span-output}) & non-member law fixed on the span & one-time build; rank condition; zero per-query \\
\bottomrule
\end{tabular}
\end{table}

\section{Notes and Provenance}
\label{sec:ch12-notes}

The linear-retrieval view, the non-member-output-on-the-span theorem, the
rank threshold and cogirth robustness, and the frequency-hiding game are the
codec paper's \cite{towell2026codec}, summarized here with the proofs left to
it. Two threads of the book meet in this chapter. The frequency-hiding is a
structural alternative to the marginal-scale homophonic lever of
\Cref{ch:marginal-scale}, achieving at build time, and at zero per-query
cost, what multiplicity achieves by spending tokens; and the opposite-pull
between the frequency game and the coincidence oracle of
\Cref{ch:compositional-scale} is the same two-scale tension the
confidentiality part drew, now appearing as a design frontier for a single
public codec. The construction is online in the sense of
\Cref{sec:strategies}, read off a linear structure rather than seed-searched,
and it is the cleanest evidence that the framework's confidentiality knobs
are not tied to any one construction strategy.
